\begin{pastexam}[label={px:ntu-108-ee-dsb-16}]{臺大電機所 108 學年度（B）第 16 題}
（複選；每題 10 分，每個選項單獨計分，答對 2 分，答錯倒扣 2 分）

Which of the following statements are true?
\begin{enumerate}[label=(\Alph*).]
\item An algorithm is $O(\log n)$ if it takes a constant time to cut the problem size by a fraction.
\item The Horner's method of evaluating the polynominal
\[ p(x) = a_0x^n + a_1x^{n-1} + a_2x^{n-2} + \cdots + a_{n-1}x + a_n \]
by
\[ p(x) = (\cdots(((a_0)\,x + a_1)\,x + a_2)\,x + \cdots + a_{n-1})\,x + a_n \]
runs in $\Theta(n^2)$ time.
\item The running time of the following recursive pseudo-code as a function of $n$ is
$\Theta(n^3)$.
\begin{minted}[fontsize=\small]{python}
def function(n):
    if n < 2: return
    else: counter = 0
    for i in range(1,8):
        function(n/2)
    for i in range(1,n*n*n):
        counter += 1
\end{minted}
\item The complexity of the following function is $O(n\log n)$.
\begin{minted}[fontsize=\small]{python}
def function(n):
    count = 0
    for i in range(n/2,n):
        j = 1
        while j + n/2 <= n:
            k = 1
            while k <= n:
                count += 1
                k *= 2
            j *= 2
    print(count)
\end{minted}
\item The complexity of the following function is $O(\log n)$.
\begin{minted}[fontsize=\small]{python}
def Function(n):
i = s = 1
while s < n:
    i += 1
    s += i
    print('*')
\end{minted}
\end{enumerate}
\end{pastexam}

\begin{pastsoln}
答案是 (A)(C)。

說明兩點。 第一，程式裡的 \texttt{n/2} 照題意當成整數除法（Python~2 的寫法；Python~3 要寫
\texttt{n//2}，否則 \texttt{range} 會出錯），不影響漸近分析。 第二，(E) 的程式除了第一行，其餘
各行原卷就沒有縮排（\texttt{i = s = 1} 和 \texttt{while} 都頂格），這裡照題意把它們都當成函式的
本體。 (B) 的 Horner 法在資結書第~4.4~節，(C) 的遞迴式在資結書第~2~章，其餘是第~1~章的程式片段
分析。

\textbf{(A) 對}。 「花常數時間把問題規模縮成原來的一個固定比例」，就是
$T(n) = T(n/b) + O(1)$（$b > 1$ 是常數；例如二分搜尋 $b = 2$）。 展開 $k$ 次得
$T(n) = T(n/b^k) + k \cdot O(1)$，到規模剩常數要 $k = \log_b n$ 次，所以 $T(n) = O(\log n)$
（也可以用主定理的情況 2，資結書第~2.4~節）。

\textbf{(B) 錯}。 Horner 法從最內層的括號開始，每一層做一次乘法、一次加法：
$r \gets a_0$，然後 $i = 1, 2, \ldots, n$ 各做 $r \gets r \cdot x + a_i$。 共 $n$ 次乘法和 $n$
次加法，$\Theta(n)$（資結書第~4.4~節、虛擬碼~4.7），不是 $\Theta(n^2)$。 $\Theta(n^2)$ 是逐項
用 $n - i$ 次乘法算 $x^{n-i}$ 的直接做法。

\textbf{(C) 對}。 \texttt{range(1,8)} 是 $i = 1, 2, \ldots, 7$，遞迴呼叫 $7$ 次，規模各是 $n/2$；
\texttt{range(1,n*n*n)} 跑 $n^3 - 1$ 次。 所以
\[
T(n) = 7\,T(n/2) + \Theta(n^3).
\]
用主定理（資結書定理~2.4）：$a = 7$、$b = 2$，$n^{\log_2 7} \approx n^{2.807}$，而
$f(n) = n^3 = \Omega(n^{\log_2 7 + \epsilon})$（取 $\epsilon = 0.19$）；正則條件
$7 f(n/2) = \frac{7}{8} n^3 \le c\,n^3$ 對 $c = 7/8 < 1$ 成立。 屬於情況 3，
$T(n) = \Theta(n^3)$。

\textbf{(D) 錯}。 三層迴圈互相獨立，次數直接相乘（資結書第~1.7~節）：
\begin{itemize}
\item 外層 \texttt{for i in range(n/2,n)}：$n - n/2 = n/2$ 次。
\item 中層 \texttt{while j + n/2 <= n}，也就是 $j \le n/2$，$j = 1, 2, 4, \ldots$：
$\lfloor \lg(n/2) \rfloor + 1$ 次。
\item 內層 \texttt{while k <= n}，$k = 1, 2, 4, \ldots$：$\lfloor \lg n \rfloor + 1$ 次。
\end{itemize}
$n$ 是 $2$ 的次方時 \texttt{count} $= \frac{n}{2} \cdot \lg n \cdot (\lg n + 1)$，例如 $n = 1024$ 得
$512 \cdot 10 \cdot 11 = 56320$。 所以複雜度是 $\Theta(n \log^2 n)$。 $n \log^2 n / (n \log n) =
\log n \to \infty$，不是 $O(n \log n)$。

\textbf{(E) 錯}。 迴圈跑了 $t$ 次後 $i = t + 1$，$s = 1 + 2 + \cdots + (t + 1) = (t+1)(t+2)/2$。
迴圈在 $s \ge n$ 時停，也就是 $(t+1)(t+2)/2 \ge n$ 的最小 $t$，大約 $t \approx \sqrt{2n}$。
所以複雜度是 $\Theta(\sqrt{n})$。 $\sqrt{n}$ 長得比 $\log n$ 快（資結書第~1.6~節），不是
$O(\log n)$。 例如 $n = 1024$ 時迴圈跑 $44$ 次，而 $\lg 1024 = 10$。
\end{pastsoln}
